\subsection[The proofs for Section 2]{The proofs for Section \ref{sec_matrix_bingham}}
\label{sec_proofs_Phi_bounds}

\begin{proof}[Proof of Theorem \ref{thm_Phi_m_bound_Bingham}]
We begin by applying to \eqref{eq_remainder_term} the triangle inequality, the upper bound from Proposition \ref{prop_zonal_bound}, and the inequality $|C_\kappa(A)| \le C_\kappa(A_+)$. Then we obtain
\begin{align}
\label{eq_Phi_m_bound1}
|\Phi_{d,p;m}(A,\Sigma)| &\le \sum_{k=m}^\infty \frac{1}{k!} \sum_{|\kappa|=k} \frac{|C_\kappa(A)|}{C_\kappa(I_d)} \Bigg(\prod_{i=1}^p \frac{\bigl(d^{1/2}/2\bigr)_{p \kappa_i}}{(d/2)_{p \kappa_i}}\Bigg)^{1/p} C_\kappa(I_d) \|\Sigma\|^{|\kappa|} \nonumber \\
&\le \sum_{k=m}^\infty \frac{\|\Sigma\|^k}{k!} \sum_{|\kappa|=k} C_\kappa(A_+) \Bigg(\prod_{i=1}^p \frac{\bigl(d^{1/2}/2\bigr)_{p \kappa_i}}{(d/2)_{p \kappa_i}}\Bigg)^{1/p}.
\end{align}
By \eqref{eq_d_kappa_ratio_ineq3},
\[
\left(\prod_{i=1}^p \frac{\bigl(d^{1/2}/2\bigr)_{p \kappa_i}}{(d/2)_{p \kappa_i}}\right)^{1/p} \le \alpha_p (|\kappa|!)^{1/2} \gamma_1^{|\kappa|} p^{|\kappa|/2} d^{-|\kappa|/2},
\]
and by inserting this bound into \eqref{eq_Phi_m_bound1}, we obtain
\begin{align}
\label{eq_Phi_m_bound2}
\Phi_{d,p;m}(A,\Sigma)| &\le \sum_{k=m}^\infty \frac{\|\Sigma\|^k}{k!} \sum_{|\kappa|=k} C_\kappa(A_+) \alpha_p (|\kappa|!)^{1/2} \gamma_1^{|\kappa|} p^{|\kappa|/2} d^{-|\kappa|/2} \nonumber \\
&= \alpha_p \sum_{k=m}^\infty \frac{\gamma_1^k p^{k/2} d^{-k/2} \|\Sigma\|^k}{(k!)^{1/2}} \sum_{|\kappa|=k} C_\kappa(A_+).
\end{align}

By \eqref{eq_trace_zonal},
$
\sum_{|\kappa|=k} C_\kappa(A_+) = (\tr A_+)^k$,
and, by assumption, $\|\Sigma\| \le \gamma_0 d^{r/2}$. Therefore, we obtain from \eqref{eq_Phi_m_bound2} the inequality
\begin{equation}
\label{eq_Phi_m_bound3}
\Phi_{d,p;m}(A,\Sigma)| \le \alpha_p R_m\bigl(\gamma_0 \gamma_1 p^{1/2} (\tr A_+) d^{-(1-r)/2}\bigr),
\end{equation}
where
\begin{equation}
\label{eq_R_m}
R_m(t) = \sum_{k=m}^\infty \frac{t^k}{(k!)^{1/2}},\qquad t > 0.
\end{equation}
By applying the ratio test, we find that the series \eqref{eq_R_m} converges for all $t$, and therefore~\eqref{eq_Phi_m_bound3} converges for all $d$. This establishes the bound \eqref{eq_Phi_m_bound}.

To obtain a closed-form upper bound for $R_m(t)$ we begin by applying Stirling's well-known inequality for the factorial function \cite[Theorem 1]{Impens}; viz., for all $k \ge 1$,
\[
k! \ge (2\pi)^{1/2} k^{(2k+1)/2} {\rm e}^{-k}.
\]
Inverting Stirling's inequality and applying the result to \eqref{eq_R_m}, we obtain
\begin{equation}
\label{eq_R_m_bound}
R_m(t) \le (2\pi)^{-1/4} \sum_{k=m}^\infty k^{-(2k+1)/4} {\rm e}^{k/2} t^k.
\end{equation}
We now apply the telescoping method to obtain an upper bound for the right-hand side of~\eqref{eq_R_m_bound}. For $j \ge m$, let
\begin{equation}
\label{eq_aj}
a_j = j^{-(2j+1)/4} {\rm e}^{j/2} t^j,
\end{equation}
then we obtain through straightforward algebraic manipulations
\begin{equation}
\label{eq_aj_ratio}
\frac{a_{j+1}}{a_j} = \bigl(1+j^{-1}\bigr)^{-j/2} \frac{j^{1/4}}{(j+1)^{3/4}} {\rm e}^{1/2} t.
\end{equation}
Since the function \smash{$j \mapsto \bigl(1+j^{-1}\bigr)^{-j/2}$}, $j \ge m$, is decreasing, then
\begin{equation}
\label{eq_c_m}
\bigl(1+j^{-1}\bigr)^{-j/2} \le \bigl(1+m^{-1}\bigr)^{-m/2},
\end{equation}
for all $j \ge m$, and by applying the latter inequality to \eqref{eq_aj_ratio}, we obtain
\[
\frac{a_{j+1}}{a_j} \le \bigl(1+m^{-1}\bigr)^{-m/2} {\rm e}^{1/2} t \frac{j^{1/4}}{(j+1)^{3/4}}
= c_m^{1/2} t \frac{j^{1/4}}{(j+1)^{3/4}},
\]
where \smash{$c_m = \bigl(1+m^{-1}\bigr)^{-m} {\rm e}$}. It follows that, for all $k \ge m$,
\begin{align}
\label{eq_a_k}
a_k = a_m \prod_{j=m}^{k-1} \frac{a_{j+1}}{a_j}
\le a_m \prod_{j=m}^{k-1} c_m^{1/2} t \frac{j^{1/4}}{(j+1)^{3/4}}= a_m c_m^{(k-m)/2} t^{k-m} \prod_{j=m}^{k-1} \frac{j^{1/4}}{(j+1)^{3/4}}.
\end{align}
Since \smash{$\prod_{j=m}^{k-1} j = (k-1)!/(m-1)!$}, then
\begin{align}
\label{eq_a_k_factor}
\prod_{j=m}^{k-1} \frac{j^{1/4}}{(j+1)^{3/4}} &= \frac{[(k-1)!/(m-1)!]^{1/4}}{[k!/m!]^{3/4}}
= \frac{m^{3/4} [(m-1)!]^{1/2}}{k^{3/4} [(k-1)!]^{1/2}}.
\end{align}
On applying to \eqref{eq_a_k} the identity \eqref{eq_a_k_factor} and the explicit expression for $a_m$, as given in \eqref{eq_aj}, we obtain
\begin{align}
\label{eq_a_k_bound}
a_k &\le m^{-(2m+1)/4} {\rm e}^{m/2} t^m \cdot c_m^{(k-m)/2} t^{k-m} \frac{m^{3/4} [(m-1)!]^{1/2}}{k^{3/4} [(k-1)!]^{1/2}} \nonumber \\
&= ({\rm e}/m c_m)^{m/2} (m!)^{1/2} \frac{c_m^{k/2} t^k}{k^{3/4} [(k-1)!]^{1/2}}.
\end{align}
Therefore, by \eqref{eq_R_m_bound}, \eqref{eq_aj}, and \eqref{eq_a_k_bound},
\begin{align}
\label{eq_sums_2_bound}
R_m(t) &\le (2\pi)^{-1/4} ({\rm e}/m c_m)^{m/2} (m!)^{1/2} \sum_{k=m}^\infty \frac{c_m^{k/2} t^k}{k^{3/4} [(k-1)!]^{1/2}} \nonumber \\
&\le (2\pi)^{-1/4} ({\rm e}/m c_m)^{m/2} (m!)^{1/2} \Bigg(\sum_{k=m}^\infty \frac{1}{k^{3/2}}\Bigg)^{1/2} \Bigg(\sum_{k=m}^\infty \frac{c_m^k t^{2k}}{(k-1)!}\Bigg)^{1/2},
\end{align}
where the latter inequality follows by the Cauchy--Schwarz inequality.

To obtain an upper bound for the first sum in \eqref{eq_sums_2_bound}, we observe from graphical considerations using Riemann sums that, for $m \ge 2$,
\begin{equation}
\label{eq_zeta_32}
\sum_{k=m}^\infty \frac{1}{k^{3/2}} \le \int_m^\infty \frac{\dd t}{(t-1)^{3/2}} = \frac{2}{(m-1)^{1/2}} \le \left(\frac{8}{m}\right)^{1/2}.
\end{equation}
As for the second sum in \eqref{eq_sums_2_bound}, we apply the inequality
$
\frac{1}{(k+m-1)!} \le \frac{1}{k! (m-1)!}
$
to obtain
\begin{align}
\label{eq_partial_exp}
\sum_{k=m}^\infty \frac{t^{2k}}{(k-1)!} &= \sum_{k=0}^\infty \frac{t^{2(k+m)}}{(k+m-1)!} \le \frac{t^{2m}}{(m-1)!} \sum_{k=0}^\infty \frac{t^{2k}}{k!}
= \frac{t^{2m}}{(m-1)!} \exp\bigl(t^2\bigr).
\end{align}

Applying \eqref{eq_zeta_32} and \eqref{eq_partial_exp} to \eqref{eq_sums_2_bound}, we obtain
\begin{align}
R_m(t) &\le (2\pi)^{-1/4} ({\rm e}/m c_m)^{m/2} (m!)^{1/2} \left(\frac{8}{m}\right)^{1/4} \left(\frac{c_m^m t^{2m}}{(m-1)!} \exp\bigl(c_m t^2\bigr)\right)^{1/2} \nonumber \\
&= (4 {\rm e}/\pi)^{1/4} ({\rm e}/m)^{(m/2)-(1/4)} t^m \exp\bigl(c_m t^2/2\bigr),\label{eq_R_m_bound2}
\end{align}
and by inserting this bound at \eqref{eq_Phi_m_bound3}, we obtain \eqref{eq_Phi_m_boundw}.

Finally, it follows from \eqref{eq_Phi_m_bound} that \smash{$\Phi_{d,p;m}(A,\Sigma) = O\bigl(d^{-(1-r)m/2}\bigr)$} as $d \to \infty$.
\end{proof}
